
\documentclass{article}
\usepackage{colortbl}
\usepackage{makecell}
\usepackage{multirow}
\usepackage{supertabular}

\begin{document}

\newcounter{utterance}

\centering \large Interaction Transcript for game `referencegame', experiment `random\_grids', episode 4 with reflexion{-}Qwen3.5{-}4B.
\vspace{24pt}

{ \footnotesize  \setcounter{utterance}{1}
\setlength{\tabcolsep}{0pt}
\begin{supertabular}{c@{$\;$}|p{.15\linewidth}@{}p{.15\linewidth}p{.15\linewidth}p{.15\linewidth}p{.15\linewidth}p{.15\linewidth}}
    \# & $\;$A & \multicolumn{4}{c}{Game Master} & $\;\:$B\\
    \hline

    \theutterance \stepcounter{utterance}  
    & & \multicolumn{4}{p{0.6\linewidth}}{
        \cellcolor[rgb]{0.9,0.9,0.9}{
            \makecell[{{p{\linewidth}}}]{
                \texttt{\tiny{[P1$\langle$GM]}}
                \texttt{You are given three grids, where each of them is 5 by 5 in size.} \\
\texttt{Grids have empty cells marked with "▢" and filled cells marked with "X".} \\
\texttt{Your task is to generate a referring expression that best describes the target grid while distinguishing it from the two other distractor grids.} \\
\texttt{The first grid is the target grid, and the following two grids are the distractors.} \\
\\ 
\texttt{Target grid:} \\
\\ 
\texttt{X X □ □ X} \\
\texttt{□ X X □ X} \\
\texttt{X □ X X □} \\
\texttt{□ X □ X X} \\
\texttt{X □ □ X X} \\
\\ 
\texttt{Distractor grid 1:} \\
\\ 
\texttt{X X □ X □} \\
\texttt{□ X X □ X} \\
\texttt{X □ □ X □} \\
\texttt{□ X □ □ X} \\
\texttt{X □ X □ X} \\
\\ 
\texttt{Distractor grid 2:} \\
\\ 
\texttt{□ X □ X □} \\
\texttt{X □ X □ X} \\
\texttt{□ X □ X □} \\
\texttt{X □ X □ X} \\
\texttt{□ X □ X □} \\
\\ 
\texttt{Instruction: Describe the target grid.} \\
\texttt{Generate the referring expression starting with the tag "Expression: " for the given target grid. Omit any other text.} \\
            }
        }
    }
    & \\ \\

    \theutterance \stepcounter{utterance}  
    & \multicolumn{4}{p{0.6\linewidth}}{
        \cellcolor[rgb]{0.9,0.9,0.9}{
            \makecell[{{p{\linewidth}}}]{
                \texttt{\tiny{[P1$\rangle$GM]}}
                \texttt{Expression: The grid with a filled cell in the top{-}left corner, a filled cell in the top{-}right corner, and a filled cell in the bottom{-}left corner, where the second row has filled cells in the second and third columns, and the fourth row has filled cells in the second and fourth columns.} \\
            }
        }
    }
    & & \\ \\

    \theutterance \stepcounter{utterance}  
    & & & \multicolumn{2}{p{0.3\linewidth}}{
        \cellcolor[rgb]{0.9,0.9,0.9}{
            \makecell[{{p{\linewidth}}}]{
                \texttt{\tiny{[GM$|$GM]}}
                \texttt{Expression: The grid with a filled cell in the top{-}left corner, a filled cell in the top{-}right corner, and a filled cell in the bottom{-}left corner, where the second row has filled cells in the second and third columns, and the fourth row has filled cells in the second and fourth columns.} \\
            }
        }
    }
    & & \\ \\

    \theutterance \stepcounter{utterance}  
    & & \multicolumn{4}{p{0.6\linewidth}}{
        \cellcolor[rgb]{0.9,0.9,0.9}{
            \makecell[{{p{\linewidth}}}]{
                \texttt{\tiny{[GM$\rangle$P2]}}
                \texttt{You are given three grids, where each of them is 5 by 5 in size.} \\
\texttt{Grids have empty cells marked with "▢" and filled cells marked with "X".} \\
\texttt{You are also given a referring expression that describes one of the given grids.} \\
\texttt{Your task is to select the grid that matches the given referring expression.} \\
\texttt{Generate only the number (in text) of the grid that matches the given expression by selecting first, second, or third.} \\
\\ 
\texttt{First:} \\
\\ 
\texttt{X X □ X □} \\
\texttt{□ X X □ X} \\
\texttt{X □ □ X □} \\
\texttt{□ X □ □ X} \\
\texttt{X □ X □ X} \\
\\ 
\texttt{Second:} \\
\\ 
\texttt{X X □ □ X} \\
\texttt{□ X X □ X} \\
\texttt{X □ X X □} \\
\texttt{□ X □ X X} \\
\texttt{X □ □ X X} \\
\\ 
\texttt{Third:} \\
\\ 
\texttt{□ X □ X □} \\
\texttt{X □ X □ X} \\
\texttt{□ X □ X □} \\
\texttt{X □ X □ X} \\
\texttt{□ X □ X □} \\
\\ 
\texttt{Expression: The grid with a filled cell in the top{-}left corner, a filled cell in the top{-}right corner, and a filled cell in the bottom{-}left corner, where the second row has filled cells in the second and third columns, and the fourth row has filled cells in the second and fourth columns.} \\
\texttt{Question: Which grid does the expression refer to?} \\
\texttt{Start with the tag "Answer: ", followed by your selection. Omit any other text.} \\
            }
        }
    }
    & \\ \\

    \theutterance \stepcounter{utterance}  
    & & & \multicolumn{4}{p{0.6\linewidth}}{
        \cellcolor[rgb]{0.9,0.9,0.9}{
            \makecell[{{p{\linewidth}}}]{
                \texttt{\tiny{[GM$\langle$P2]}}
                \texttt{Answer: First} \\
            }
        }
    }
     \\ \\

    \theutterance \stepcounter{utterance}  
    & & & \multicolumn{2}{p{0.3\linewidth}}{
        \cellcolor[rgb]{0.9,0.9,0.9}{
            \makecell[{{p{\linewidth}}}]{
                \texttt{\tiny{[GM$|$GM]}}
                \texttt{first} \\
            }
        }
    }
    & & \\ \\

    \theutterance \stepcounter{utterance}  
    & \multicolumn{4}{p{0.6\linewidth}}{
        \cellcolor[rgb]{0.95,0.95,0.95}{
            \makecell[{{p{\linewidth}}}]{
                \texttt{\tiny{[P2$\rightarrow$P1]}}
                \texttt{Post{-}episode Reflexion (not sent to GM):} \\
\texttt{New plan: My previous description was too verbose and likely failed because it listed every single cell's state rather than highlighting the unique distinguishing features of the target grid compared to the distractors. I should have focused on the specific pattern that makes the target grid unique, such as the filled cell in the bottom{-}left corner (row 5, col 1) which is empty in both distractors, or the specific arrangement of the top row. Next time, I will generate a concise expression that identifies the target grid by its most distinctive feature (e.g., "The grid with filled cells at the top{-}left, top{-}right, and bottom{-}left corners, and a filled cell in the fourth row, fourth column") without listing every coordinate, ensuring the description clearly separates it from the distractors.} \\
            }
        }
    }
    & & \\ \\

    \theutterance \stepcounter{utterance}  
    & & & \multicolumn{4}{p{0.6\linewidth}}{
        \cellcolor[rgb]{0.95,0.95,0.95}{
            \makecell[{{p{\linewidth}}}]{
                \texttt{\tiny{[P1$\rightarrow$P2]}}
                \texttt{Post{-}episode Reflexion (not sent to GM):} \\
\texttt{New plan: I incorrectly identified the grid by relying on a partial match of the top row and bottom{-}left corner, while ignoring the specific constraints for the second and fourth rows. The correct strategy is to systematically verify every condition in the expression against each grid row by row: first check the top{-}left and top{-}right corners, then confirm the filled cells in the second row (columns 2 and 3), and finally verify the filled cells in the fourth row (columns 2 and 4). Only if all these specific coordinates match exactly should I select that grid.} \\
            }
        }
    }
     \\ \\

\end{supertabular}
}

\end{document}
